% Generated by scripts/assemble.py; corrections belong in source/corrections.json.
\examhead{190}{西安电子科技大学 2024 年数学分析试卷}{179}
% p180-b19; source PDF page 180
\par\Needspace{5\baselineskip}
\blockanchor{p180-b19}
\majorquestion{一}计算题. 每题 7 分, 共 42 分.\par
% p180-b20; source PDF page 180
\blockanchor{p180-b20}
\question{1}求极限 $\lim_{n\to \infty}\sqrt[n]{1 + a^n +\left(\frac{a^2}{2}\right)^n}$$(a > 0)$.\par
% p180-b21; source PDF page 180
\blockanchor{p180-b21}
\question{2}求不定积分 $\int e^{-\sin x} \frac{\sin 2x}{\sin^4\left(\frac{\pi}{4} - \frac{x}{2}\right)} \mathrm{d}x.$\par
% p180-b22; source PDF page 180
\blockanchor{p180-b22}
\question{3}求极限 $\lim_{x\to 0}\frac{6\left(\sqrt[3]{2 - \cos x} - 1\right) + \ln(1 - x^2)}{x^4}$.\par
% p180-b23; source PDF page 180
\blockanchor{p180-b23}
\question{4}求重积分 $\iint_{D} e^{\frac{y}{x + y}} \, \mathrm{d}x \, \mathrm{d}y$, 其中 $D$ 是 $x + y = 1, x = 0, y = 0$ 所围成的区域.\par
% p180-b24; source PDF page 180
\blockanchor{p180-b24}
\question{5}求曲线积分 $\oint_{L} \frac{x \, \mathrm{d}y - y \, \mathrm{d}x}{4x^{2} + y^{2}}$, 其中 $L$ 是圆周 $(x - 1)^{2} + y^{2} = R^{2}$, 取逆时针方向.\par
% p180-b25; source PDF page 180
\par\Needspace{6\baselineskip}
\blockanchor{p180-b25}
\question{6}设 $\Sigma$ 是 $z = 2 - x^2 - y^2$ 介于 $z = 1$ 和 $z = 2$ 之间的部分, 取上侧, 求曲面积分\par
\nopagebreak[4]
% p180-b26; source PDF page 180
\blockanchor{p180-b26}
\displayformula{\iint_ {\Sigma} \left(x ^ {3} z + x\right) \mathrm{d} y \mathrm{d} z - x ^ {2} y z \mathrm{d} z \mathrm{d} x - x ^ {2} z ^ {2} \mathrm{d} x \mathrm{d} y.}
% p181-b00; source PDF page 181
\par\Needspace{5\baselineskip}
\blockanchor{p181-b00}
\majorquestion{二}判断下列命题是否正确, 若正确请给出证明, 若错误则给出反例. 每题 6 分, 共 24 分.\par
% p181-b01; source PDF page 181
\blockanchor{p181-b01}
\question{1}若数列 $\{a_{n}\}$ 以0为极限, 数列 $\{b_{n}\}$ 收敛, 则级数 $\sum_{n=1}^{\infty}(a_{n}b_{n})$ 收敛.\par
% p181-b02; source PDF page 181
\blockanchor{p181-b02}
\question{2}设 $f: \mathbb{R} \to \mathbb{R}$ 是连续映射, $f_n = f \circ f \circ \cdots \circ f$ 为 $f$ 的 $n$ 次复合, 若存在 $x_0 \in \mathbb{R}$ 使得数列 $\{f_n(x_0)\}$ 收敛到 $a$, 则 $a$ 是 $f$ 的不动点, 即 $f(a) = a$.\par
% p181-b03; source PDF page 181
\blockanchor{p181-b03}
\question{3}设 $f(x), f_n(x) (n = 1,2,\dots)$ 均在 $[0,1]$ 上连续，且 $\{f_n(x)\}$ 在 $[0,1]$ 上收敛于 $f(x)$，则 $\{f_n(x)\}$ 在 $[0,1]$ 上一致收敛于 $f(x)$.\par
% p181-b04; source PDF page 181
\blockanchor{p181-b04}
\question{4}设 $D$ 是凸区域, 若在 $D$ 内满足 $f_{x}(x,y) = 0$, 则 $f(x,y)$ 在 $D$ 内的取值与 $x$ 无关.\par
% p181-b05; source PDF page 181
\par\Needspace{5\baselineskip}
\blockanchor{p181-b05}
\majorquestion{三}(12 分) 讨论 $a, b, c$ 取何值时, $\sum_{n=3}^{\infty} \frac{a^n}{n^b (\ln n)^c}$ 绝对收敛, 条件收敛, 发散.\par
% p181-b06; source PDF page 181
\par\Needspace{5\baselineskip}
\blockanchor{p181-b06}
\majorquestion{四}(12 分) 二元函数 $f(x, y)$ 在点 (1,0) 的某邻域有一阶连续偏导数, 且 $f(1,0) = 0, f_y(1,0) \neq 0$, 证明: 方程 $f\left(x, \int_0^t \cos u \, \mathrm{d}u\right) = 0$ 在点 (1,0) 的某邻域内可确定隐函数 $t = \varphi(x)$, 并求 $\varphi'(1)$.\par
% p181-b07; source PDF page 181
\par\Needspace{5\baselineskip}
\blockanchor{p181-b07}
\majorquestion{五}(15 分) 设 $f(x, y, z) = \ln x + 2\ln y + \ln z$, 求 $f(x, y, z)$ 在 $x^2 + y^2 + z^2 = 4r^2 (x, y, z > 0)$ 下的极大值, 其中 $r$ 为正常数. 并证明对任意的正实数 $a, b, c$, 满足 $ab^2c \leq 4\left(\frac{a + b + c}{4}\right)^4$.\par
% p181-b08; source PDF page 181
\par\Needspace{5\baselineskip}
\blockanchor{p181-b08}
\majorquestion{六}(15 分) 设 $f(x) = \sum_{n=0}^{\infty} \frac{1}{2^n + x}$.\par
% p181-b09; source PDF page 181
\blockanchor{p181-b09}
(1) 证明: $f(x)$ 在 $[0, +\infty)$ 上可导且一致连续.\par
% p181-b10; source PDF page 181
\blockanchor{p181-b10}
(2) 证明: 反常积分 $\int_{0}^{+\infty} f(x) \, \mathrm{d}x$ 发散.\par
% p181-b11; source PDF page 181
\par\Needspace{5\baselineskip}
\blockanchor{p181-b11}
\majorquestion{七}(15 分) 设 $f(x)$ 在 $[0,2]$ 上连续, $h \in \left(0, \frac{1}{2}\right)$, 记 $f_h(x) = \frac{1}{h^2} \int_0^h \mathrm{d}u_2 \int_0^h f(x + u_1 + u_2) \mathrm{d}u_1$.\par
% p181-b12; source PDF page 181
\blockanchor{p181-b12}
(1) 证明: $f_{h}(x)$ 在 $[0,1]$ 上二阶可导, 并求 $f_{h}^{\prime \prime}(x)$.\par
% p181-b13; source PDF page 181
\blockanchor{p181-b13}
(2) 若 $f(x)$ 在 $x_0 \in [0,1]$ 处存在二阶导数, 证明: $\lim_{h \to 0^+} f_h''(x_0) = f''(x_0)$.\par
% p181-b14; source PDF page 181
\par\Needspace{5\baselineskip}
\blockanchor{p181-b14}
\majorquestion{八}(15 分) 设 $f(x)$ 在 $[a, b]$ 上黎曼可积.\par
% p181-b15; source PDF page 181
\blockanchor{p181-b15}
(1) 证明: 对任意的 $\varepsilon > 0$, 存在区间 $[c, d] \subset [a, b]$, 使得 $f(x)$ 在 $[c, d]$ 振幅 $\omega_f < \varepsilon$.\par
% p181-b16; source PDF page 181
\blockanchor{p181-b16}
(2) 证明: $f(x)$ 的连续点在 $[a, b]$ 上稠密, 即对任意的 $[\alpha, \beta] \subset [a, b]$, $f(x)$ 在 $[\alpha, \beta]$ 内部有连续点.\par
% p181-b17; source PDF page 181
\blockanchor{p181-b17}
(3) 若 $f(x) \geq 0$, 证明: $\int_{a}^{b} f(x) \, \mathrm{d}x = 0$ 当且仅当 $f(x)$ 在其连续点恒为 0.\par
